\section{Flüssigmetall im Magnetfeld: MHD in FluidX3D}\label{sec:mhd}

Im Abschnitt über flüssiges Blei (Abschnitt~\ref{sec:filament}) wurde der Druckverlust im Magnetfeld nur mit
einer Faustformel abgeschätzt: $\mathrm{d}p/\mathrm{d}x = \sigma\,v\,B^2 k_p$, mit $k_p \approx 1/\mathrm{Ha}$ für
isolierende Wände. Der Nutzer fragte, ob sich die Strömung im Magnetfeld auch dynamisch rechnen lässt, als eine Art
CFD-Rechnung, und ob man das mit FluidX3D koppeln kann. FluidX3D kennt von sich aus keine Magnetfelder. Dieser
Abschnitt beschreibt, wie die nötige Kraft in zwei Stufen eingebaut und an Fällen mit bekannter Lösung geprüft wurde.

\paragraph{Was passiert physikalisch?} Bewegt sich ein elektrisch leitendes Metall quer durch ein Magnetfeld,
entsteht darin eine elektrische Spannung, wie in einem Dynamo. Sie treibt Ströme $\mathbf{j}$, und diese Ströme
erfahren im Feld wieder eine Kraft, die \emph{Lorentzkraft} $\mathbf{j}\times\mathbf{B}$. Sie bremst die Bewegung
quer zum Feld. Bei kleiner magnetischer Reynoldszahl, wie bei Blei in einer Spule, verändern die Ströme das äußere
Feld kaum (\emph{induktionsfreie} Näherung). Dann gilt das Ohmsche Gesetz für bewegte Leiter und die Erhaltung der
Ladung:
\[
\mathbf{j} = \sigma\left(-\nabla\varphi + \mathbf{u}\times\mathbf{B}\right),\qquad \nabla\cdot\mathbf{j} = 0
\quad\Rightarrow\quad \nabla^2\varphi = \nabla\cdot(\mathbf{u}\times\mathbf{B}),\qquad
\mathbf{f} = \frac{\mathbf{j}\times\mathbf{B}}{\rho}.
\]
$\sigma$ ist die elektrische Leitfähigkeit, $\varphi$ das elektrische Potential. Wie stark das Feld gegenüber der
Zähigkeit wirkt, misst die \emph{Hartmann-Zahl} $\mathrm{Ha} = B\,a\sqrt{\sigma/(\rho\nu)}$ ($a$ halbe Kanalweite in
Feldrichtung). An den Wänden senkrecht zum Feld (\emph{Hartmann-Wände}) bilden sich dünne Grenzschichten der Dicke
$a/\mathrm{Ha}$, an den Wänden parallel zum Feld (\emph{Seitenwände}) dickere Schichten der Dicke $a/\sqrt{\mathrm{Ha}}$
\cite{huntshercliff1971}. Ob sich die Ströme durch die Wände oder nur durch diese Schichten schließen, entscheidet über
den Druckverlust.

\paragraph{Begriffe in diesem Abschnitt.}
\begin{description}\itemsep2pt
\item[Hartmann-Wände] die beiden Kanalwände, auf die das Magnetfeld senkrecht trifft. Zeigt das Feld nach oben, sind
  es Boden und Decke des Kanals. Benannt nach dem dänischen Physiker Julius Hartmann, der diese Strömung in den 1930er-Jahren untersuchte.
\item[Seitenwände] die beiden anderen Wände, an denen das Feld entlangläuft (links und rechts).
\item[isolierende Wand] eine Wand, durch die kein elektrischer Strom fließen kann, etwa mit einer Keramikschicht.
  Die im Metall erzeugten Ströme müssen sich dann innerhalb des Metalls schließen, in den dünnen Schichten an der Wand.
\item[leitende Wand] eine Wand, die den Strom gut leitet, etwa aus Kupfer oder dickem Stahl. Die Ströme können sich
  dann über die Wand schließen wie über einen Draht; der Stromkreis hat wenig Widerstand, und die Bremskraft wird groß.
  \emph{Ideal leitend} heißt: ohne jeden Widerstand, ein Grenzfall für die Rechnung.
\item[leitende Hartmann-Wände] Boden und Decke leiten, die Seitenwände isolieren. Das ist der Fall von Hunt (1965), für
  den es eine genaue Lösung gibt.
\item[Hartmann-Zahl Ha] ein Maß dafür, wie stark das Magnetfeld gegenüber der Zähigkeit des Metalls wirkt. Ha = 0 heißt
  kein Feld; bei Ha = 100 beherrscht das Feld die Strömung fast vollständig.
\item[Hartmann-Schicht] die dünne Schicht an den Hartmann-Wänden, in der die Geschwindigkeit von null auf den Wert in
  der Mitte steigt. Ihre Dicke ist die halbe Kanalweite geteilt durch Ha, bei Ha = 100 also ein Hundertstel.
\item[Seitenschicht] die entsprechende Schicht an den Seitenwänden. Sie ist dicker, halbe Kanalweite geteilt durch
  $\sqrt{\mathrm{Ha}}$.
\item[Potential $\varphi$] die elektrische Spannung an jeder Stelle im Metall, gemessen gegen einen festen Punkt. Ihr
  Gefälle lenkt die Ströme so, dass sie sich schließen.
\item[Jacobi-Schritt] ein einfaches Rechenverfahren für $\varphi$: Jede Zelle nimmt immer wieder den Mittelwert ihrer
  Nachbarn an, korrigiert um die Quelle, bis sich nichts mehr ändert.
\item[Druckverlustbeiwert $k_p$] gibt an, welcher Bruchteil der größtmöglichen magnetischen Bremskraft tatsächlich wirkt
  (1 = volle Bremsung, klein = die Ströme finden einen schwierigen Rückweg).
\end{description}

\subsection*{Stufe 1: Lorentzkraft ohne Potential}
Zwischen zwei unendlich ausgedehnten Platten verschwindet $\nabla\varphi$, und es bleibt
$\mathbf{f} = -K\,\mathbf{u}_\perp$ mit $K = \sigma B^2/\rho$. In die eigene Kopie von FluidX3D wurde diese Kraft in den
Hauptkernel eingebaut, halbimplizit zur Guo-Kraft: $\mathbf{u}_\perp = (\mathbf{u}^*_\perp + \mathbf{F}_\perp/2\rho)/(1+K/2)$.
So bleibt die Rechnung auch bei großem $K$ stabil. Die Änderungen liegen als Diffs im Ordner \texttt{fluidx3d/}
(FluidX3D-Lizenz \cite{fluidx3d}). Geprüft wurde an der Hartmann-Strömung, für die es eine exakte Lösung gibt:
$u(y) = f_x/K\,(1-\cosh(\mathrm{Ha}\,y/a)/\cosh\mathrm{Ha})$.

\bild{mhd_hartmann.png}{1.0}{Stufe 1: Hartmann-Strömung zwischen zwei Platten. Links FluidX3D (Punkte) gegen die exakte
Lösung (Linien) für Ha = 0 bis 30. Rechts der Fehler bei Ha = 30 über der Gitterfeinheit.}
\lesehilfe{Links ist jede Farbe eine Feldstärke. Waagrecht steht die Lage zwischen den beiden Platten ($-1$ und $+1$
sind die Wände), senkrecht die Geschwindigkeit, bezogen auf ihren Höchstwert. Ohne Feld (grau) ergibt sich die
bekannte Parabel. Je stärker das Feld, desto flacher wird das Profil in der Mitte und desto steiler fällt es an den
Wänden ab: Das Feld bremst die schnelle Mitte, nur in dünnen Schichten an der Wand bleibt ein Gefälle. Die Punkte
(FluidX3D) liegen auf den Linien (exakte Lösung). Rechts sieht man, wie der Fehler schrumpft, wenn man das Gitter
verfeinert. Die Beschriftung nennt, wie viele Rechenzellen in der dünnen Wandschicht liegen. Selbst mit nur einer
Zelle bleibt der Fehler unter 1\,\%; bei feinerem Gitter fällt er etwa wie die gestrichelte Linie (zweite Ordnung).}

Ohne Feld beträgt die Abweichung 0,03\,\%, bei Ha = 5 0,04\,\% und bei Ha = 30 mit 32 Zellen, also gut einer Zelle
in der Hartmann-Schicht, 0,5\,\%. Bei feinerem Gitter fällt der Fehler etwa mit der zweiten Ordnung (3,0 / 1,5 / 0,48
/ 0,13\,\% für $a$ = 8 / 16 / 32 / 64 Zellen).

\paragraph{Korrektur.} Die zuerst ermittelten Werte waren deutlich schlechter (4,3\,\% bei Ha = 30 und 32 Zellen,
68\,\% bei 8 Zellen) und standen so schon in einer früheren Fassung. Sie waren ein Fehler der Auswertung, nicht der
Rechnung: Ohne den Schalter \texttt{UPDATE\_FIELDS} berechnet FluidX3D die ausgegebene Geschwindigkeit in einem
eigenen Kernel nach, und dort fehlte die halbe Lorentzkraft. Je größer $K$, desto größer der Fehler. Aufgefallen ist
das erst, als der Hartmann-Test mit dem Programmstand von Stufe~2 wiederholt wurde, in dem \texttt{UPDATE\_FIELDS} an
ist. Die Regel „mindestens zwei Zellen in der Hartmann-Schicht“ war damit zu streng; etwa eine Zelle genügt für
1\,\% Genauigkeit.

\subsection*{Stufe 2: Kanäle mit Seitenwänden}
In einem echten Kanal fließen die Ströme im Kreis. In der Mitte laufen sie quer zur Strömung, zurück müssen sie über
die Wände oder durch die dünnen Schichten an den Wänden. Diesen Rückweg beschreibt das Potential $\varphi$. Ohne
$\varphi$ würde Stufe~1 die Bremskraft in einem Kanal mit isolierenden Wänden um ein Vielfaches überschätzen.

\paragraph{Umsetzung.} Ein neuer Schalter \texttt{MHD\_POT} legt zwei Puffer für $\varphi$ an. Vor jedem
Zeitschritt laufen auf der GPU einige Jacobi-Schritte für $\nabla\cdot\mathbf{j} = 0$, hier 20 Doppelschritte. Die
Gleichung ist als Bilanz der Ströme über die sechs Flächen jeder Zelle geschrieben. So bleibt die Ladung auch
diskret erhalten. Auf einer isolierenden Wand ist der Strom durch die Wandfläche null. Eine ideal leitende Wand hält
$\varphi = 0$ auf ihrer Oberfläche, die eine halbe Zelle vor dem Wandknoten liegt. Im Hauptkernel kommt die Kraft
$K\rho\,(\mathbf{E}\times\hat{\mathbf{B}})$ mit $\mathbf{E} = -\nabla\varphi$ aus denselben Flächenwerten hinzu.

\paragraph{Prüfung.} Ein Vergleich nur mit sich selbst wäre wertlos. Deshalb entstand ein eigener 2D-Löser für die
ausgebildete Kanalströmung, der einen anderen Weg geht: Er rechnet nicht mit $\varphi$, sondern mit dem induzierten
Magnetfeld $b$ (Shercliff \cite{shercliff1953}, Hunt \cite{hunt1965}):
\[
\nabla^2 u + \mathrm{Ha}\,\frac{\partial b}{\partial y} = -1,\qquad \nabla^2 b + \mathrm{Ha}\,\frac{\partial u}{\partial y} = 0 .
\]
Er reproduziert drei bekannte Grenzfälle: den quadratischen Kanal ohne Feld ($\bar u = 0{,}1406$, entsprechend
$f\,\mathit{Re} = 14{,}23$), bei breitem Kanal mit leitenden Wänden die exakte Hartmann-Lösung und bei breitem Kanal mit
isolierenden Wänden die Kerngeschwindigkeit $1/\mathrm{Ha}$. Dabei fiel ein Fehler in seiner eigenen Randbehandlung
auf, der behoben wurde.

Geprüft wurden zwei klassische Fälle im quadratischen Kanal. Bei \emph{Shercliff} sind alle Wände isolierend. Bei
\emph{Hunt} sind die Hartmann-Wände ideal leitend und die Seitenwände isolierend; dort fließt der Strom durch die
Wand zurück, und an den Seitenwänden entstehen schnelle Strahlen.

\paragraph{Ein Fehler, den erst die Prüfung zeigte.} Im ersten Durchlauf wuchs der Fehler bei isolierenden Wänden
stark mit Ha: 12\,\% bei Ha = 15, 36\,\% bei Ha = 30. Bei leitenden Wänden war alles gut. Die Ursache: Im Kern eines
isolierten Kanals heben sich zwei große Kräfte fast auf, die Bremskraft $-K\mathbf{u}_\perp$ und die Kraft aus dem
Potential. Übrig bleibt nur ein Teil von etwa $1/\mathrm{Ha}$. Jeder kleine Widerspruch zwischen den beiden Teilen wird
deshalb mit Ha verstärkt. Im Kernel wurde die Potentialkraft erst nach dem halbimpliziten Teil addiert, sodass diesem
ein Teil der Kraft fehlte. Nach dem Tausch der Reihenfolge, und mit derselben gespeicherten Geschwindigkeit für
$\mathbf{u}\times\mathbf{B}$ wie im Jacobi-Schritt, sank der Fehler bei Ha = 15 von 12\,\% auf 0,45\,\%.

\begin{table}[H]\centering\small
\begin{tabular}{lrrrr}\toprule
Fall & Ha & $a$ [Zellen] & Fehler Profil & Fehler Durchsatz\\\midrule
ohne Feld & 0 & 32 & 0,01\,\% & $-0{,}03$\,\%\\
Shercliff (isolierend) & 5 / 10 / 15 & 32 & 0,09 / 0,23 / 0,45\,\% & $-0{,}15$ / $-0{,}30$ / $-0{,}37$\,\%\\
Shercliff & 30 & 64 & 0,49\,\% & $-0{,}19$\,\%\\
Hunt (leitende Hartmann-Wände) & 5 / 10 / 15 & 32 & 0,08 / 0,42 / 0,86\,\% & $-0{,}09$ / $-0{,}30$ / $-0{,}68$\,\%\\
Hunt & 30 & 64 & 0,79\,\% & $-0{,}78$\,\%\\
Gitterstudie Shercliff & 10 & 16 / 32 / 64 & 0,75 / 0,23 / 0,09\,\% & \\\bottomrule
\end{tabular}\caption{Stufe 2: FluidX3D gegen den unabhängigen 2D-Löser (RMS-Fehler des Profils bezogen auf den
Höchstwert; Abweichung des mittleren Durchsatzes).}\end{table}

\bild{mhd_kanal.png}{1.0}{Stufe 2: quadratischer Kanal, B senkrecht nach oben. Oben links und Mitte isolierende Wände
(Shercliff), oben rechts leitende Hartmann-Wände (Hunt). Unten links und Mitte der ganze Querschnitt bei Hunt, Ha = 30,
FluidX3D und Referenz. Unten rechts die Gitterstudie.}
\lesehilfe{Man blickt in einen quadratischen Kanal hinein; das Metall fließt auf den Betrachter zu, das Magnetfeld
zeigt nach oben. Oben links das Geschwindigkeitsprofil von unten nach oben, also entlang des Feldes, oben in der
Mitte von links nach rechts, quer zum Feld. Die Punkte sind FluidX3D, die Linien der unabhängige Vergleichslöser; die
Prozentzahlen sind die Abweichungen. Mit isolierenden Wänden wird das Profil mit stärkerem Feld in Feldrichtung
kastenförmig flach (dünne Schichten oben und unten). Quer dazu bleibt es runder, weil die Schichten an den Seiten
dicker sind. Oben rechts sind die Wände oben und unten leitend: Jetzt bremst das Feld die Mitte stark, und das Metall
weicht in zwei schnelle Strahlen an den Seitenwänden aus. Die beiden Farbbilder unten zeigen genau das im ganzen
Querschnitt (gelb schnell, dunkel langsam); links FluidX3D, in der Mitte die Referenz. Sie sind kaum zu
unterscheiden. Unten rechts fällt der Fehler, wenn man das Gitter verfeinert.}

Alle zwölf Fälle stimmen auf besser als 1\,\% mit dem Vergleichslöser überein. Die Gitterstudie zeigt eine Ordnung
von 1,4 bis 1,7, also etwas schlechter als bei Stufe~1. Grund ist vermutlich die Behandlung der Wandflächen; das
ist noch nicht untersucht.

\paragraph{Prüfung der früheren Faustformel.} Aus dem Vergleichslöser folgt der Druckverlustbeiwert
$k_p = \mathrm{d}p/\mathrm{d}x\,/\,(\sigma\bar u B^2) = 1/(\mathrm{Ha}^2\,\bar U)$ im quadratischen Kanal:
\begin{table}[H]\centering\small
\begin{tabular}{rrrr}\toprule
Ha & $1/\mathrm{Ha}$ & $k_p$ isolierend & $k_p$ leitende Hartmann-Wände\\\midrule
5 & 0,200 & 0,407 & 0,925\\
10 & 0,100 & 0,153 & 0,634\\
20 & 0,050 & 0,065 & 0,475\\
50 & 0,020 & 0,023 & 0,335\\
100 & 0,010 & 0,011 & 0,256\\\bottomrule
\end{tabular}\caption{Druckverlustbeiwert aus dem 2D-Löser (Gitter mit mindestens drei Zellen je Hartmann-Schicht).}\end{table}
Mit isolierenden Wänden nähert sich $k_p$ tatsächlich $1/\mathrm{Ha}$; bei Ha = 100 liegt er noch 10\,\% darüber. Die
Faustformel war also für große Ha richtig und liegt bei kleinen Ha zu niedrig. Leiten die Hartmann-Wände, fällt
$k_p$ nur langsam, und der Druckverlust ist bei Ha = 100 etwa 25-mal so groß wie mit isolierenden Wänden.

\paragraph{Grenzen und nächste Schritte.}
\begin{itemize}\itemsep1pt
\item \textbf{Reale Hartmann-Zahlen sind viel größer.} Bleilithium in den \SI{10}{mm}-Kanälen bei \SI{6}{T} ergibt
  grob Ha $\approx 600$, in Blanket-Kanälen von einigen Zentimetern einige Tausend (Stoffwerte noch nicht geprüft).
  Mit einer Zelle je Hartmann-Schicht bräuchte schon Ha = 600 etwa 600 Zellen je halber Kanalweite. Der
  Speicher reicht dafür, aber der einfache Jacobi-Löser konvergiert auf so feinen Gittern zu langsam. Nötig wären ein
  Mehrgitterverfahren oder ein zur Wand verdichtetes Gitter; bis dahin gelten gerechnete Werte bis Ha $\approx 30$
  und darüber die geprüften Skalierungen.
\item \textbf{Endlich leitende Wände} (Stahl, Wandleitfähigkeit $c$) fehlen noch; bisher gibt es nur ideal isolierend
  und ideal leitend. Für Stahlkanäle ist der Fall dazwischen der wichtigste.
\item \textbf{Nur eine GPU}, und das Feld ist bisher überall gleich. Das Feld der Spulen ($\mathbf{B}(\mathbf{x})$
  aus Biot-Savart) wäre der Schritt zur Kopplung mit dem Stellarator.
\end{itemize}
